% Part: methods
% Chapter: proofs
% Section: cant-do-it

\documentclass[../../../include/open-logic-section]{subfiles}

\begin{document}

\olfileid[tr]{mth}{prf}{cnt}

\olsection{Bunu Yapamıyorum!{}}

Hepimiz pes edecekmişiz gibi hissettiğimiz bir noktaya geliriz. Ama
bunu \emph{yapabilirsiniz}. Öğretim elemanınız ve asistanınız kadar
sınıf arkadaşlarınız da size yardım edebilir. Onlardan yardım isteyin!{}
İşte bir bunalımdan kaçınmanıza yardımcı olacak ve pes etmek
istediğinizde ne yapacağınızı gösterecek birkaç öneri.

Problemleri başarıyla çözebilmek için şunları yapın:
\begin{enumerate}
\item \emph{Mümkün olduğunca erken başlayın.} Dönem boyunca işlerimiz
  artar ve çoğumuz erteleme eğilimiyle mücadele ederiz; yapabileceğiniz
  en iyi şeylerden biri ödevlerinize erkenden başlamaktır. Böylece
  takılırsanız bir çözüm aramak için vaktiniz olur (ağlamak dışında).
\item \emph{Sınıf arkadaşlarınızla konuşun}. Yalnız değilsiniz.
  Sınıftaki başkaları da zorlanıyor olabilir---ama farklı konularda
  zorlanabilirler. Problemi akranlarınızla konuşmak, size bir ilerlemeye
  yol açabilecek farklı bir bakış açısı verebilir. Elbette onların
  çözümünü yalnızca kopyalamayın: bir ipucu isteyin ya da nerede
  takıldığınızı açıklayıp sonraki adımı sorun. Çözdüğünüzde de karşılığını
  verin. Başka birine yardım etmek ve bir şeyleri açıklamak sizin de
  daha iyi anlamanıza yardımcı olur.
\item \emph{Yardım isteyin.} Yararlanabileceğiniz birçok kaynak
  vardır---öğretim elemanınız ve asistanınız size yardım etmek için
  oradadır ve başarılı olmanızı \emph{isterler}. Bir problemi çözümleme
  sürecinde size yardım edebilmeli ve sürecin neresinde zorlandığınızı
  belirleyebilmelidirler.
\item \emph{Ara verin.} Takıldıysanız bunun nedeni probleme çok uzun
  süredir bakıyor olmanız \emph{olabilir}. Kısa bir ara verin, bir fincan
  çay için ya da bir süre başka bir problem üzerinde çalışın; sonra
  probleme dinlenmiş bir zihinle dönün. Bir gece uyuyup yeniden düşünün.
\end{enumerate}

Bu stratejilerin kanıt üzerinde çok önceden çalışmaya başlamanızı
gerektirdiğine dikkat ettiniz mi? Kanıta teslim edilmesi gereken günden
önceki gece saat 2'de başladıysanız bunlar pek yardımcı olmayabilir.

Bu, fazlasıyla kasvetli gelebilir; ama bir kanıtı çözmek, sonunda
karşılığını veren bir uğraştır. Bazı insanlar bunu meslek olarak
yapıyor---demek ki bundan hoşlanılacak bir şey olmalı. Hemen her şey
gibi, problem çözmek ve kanıt yapmak da alıştırma gerektirir. Bunu kolay
bulan sınıf arkadaşlarınız olabilir: muhtemelen yalnızca daha önce çok
alıştırma yapmışlardır. Çok kolay pes etmemeye çalışın.

Belirli bir problemde vaktiniz (ya da sabrınız) tükenirse sorun değil.
Bu, aptal olduğunuz ya da onu hiçbir zaman anlayamayacağınız anlamına
gelmez. Nasıl yapıldığını (öğretim elemanınızdan ya da başka bir
öğrenciden) öğrenin ve nerede hata yaptığınızı ya da takıldığınızı
belirleyin; böylece benzer bir sorunla bir daha karşılaştığınızda bundan
kaçınabilirsiniz. Sonra çözüme bakmadan yapmayı deneyin. Bir dahaki
sefere de daha erken başlayın (ve yardım isteyin).

\end{document}
